\begin{frame}<1-4>[t]{课堂巩固}
\begin{campuscanvas}
% Source slide 21, objects 8; visible 2-4
\only<2-4|handout:1>{\node[anchor=north west,inner sep=0,text=black] at (70.000,155.000) {\begin{minipage}[t]{142.500mm}\raggedright\fontsize{12.5}{18.75}\selectfont \textbf{1.} 下面给出的四类对象中，能构成集合的是（\quad）。\par A．某班视力较好的同学\qquad B．某小区长寿的人\par C．$\pi$ 的近似值\qquad\qquad D．方程 $x^2=1$ 的实数根\end{minipage}};}
% Source slide 21, objects 10; visible 3-4
\only<3-4|handout:1>{\node[anchor=north west,inner sep=0,text=black] at (70.000,355.000) {\begin{minipage}[t]{142.500mm}\raggedright\fontsize{14}{21.00}\selectfont \red{\textbf{答案：D}}\end{minipage}};}
% Source slide 21, objects 9; visible 4
\only<4|handout:1>{\node[anchor=north west,inner sep=0,text=black] at (70.000,445.000) {\begin{minipage}[t]{142.500mm}\raggedright\fontsize{11.5}{17.25}\selectfont \red{解析：}\par A．“视力较好”不确定，不能构成集合。\par B．“长寿”不确定，不能构成集合。\par C．未给出精确度，“$\pi$ 的近似值”不确定，不能构成集合。\par D．方程 $x^2=1$ 的实数根是 $-1$ 和 $1$，明确可靠，能构成集合。故选 D。\end{minipage}};}
\end{campuscanvas}
\end{frame}
